\section{Other clocks and spatial startup profiles}\label{bcs:clock-options}
The compact construction is now complete. We record optional root clocks and startup profiles. The logarithmic estimate in Lemma~\ref{bcs:clock} gains a factor $(1+t)^{-1}$ with each time derivative; the root-clock estimate below gives the common exponent $\beta$ at every fixed mixed order. The final profile proposition acts on a fixed spatial field $Z$, whose flow is autonomous, rather than on a time-dependent processor. Its force and trajectory identities also apply to the distinct startup profile $\sigma(t)$ used for the spatial suspension in Section~\ref{bcs:expanded-section}.

\begin{proposition}[The root clocks]\label{bcs:power-clock}
Under the hypotheses of Lemma~\ref{bcs:clock}, set $g(t)=(1+t)^q-1$, where $q=1/3$ or $q=1/4$, and put $\beta=1-q$. All mixed derivatives of the velocity and force in \eqref{bcs:clock-force} are bounded, and for every fixed $j,\alpha$,
\[
 \|\partial_t^jD_x^\alpha u(t)\|_\infty+
 \|\partial_t^jD_x^\alpha f(t)\|_\infty
 \le C_{j,\alpha}(1+t)^{-\beta}.
\]
In particular each is square-integrable in time in spatial supremum norm. The cube-root exponent is $2/3$ and the fourth-root exponent is $3/4$.
\end{proposition}
\begin{proof}
Every positive-order derivative $g^{(k)}$ is bounded and is $O_k((1+t)^{-\beta})$. Each differentiated term in $g'W(g)$ contains at least one such factor and a bounded derivative of $W$. One factor supplies the decay and the others are bounded. Applying this statement to $u_t-\nu\Delta u$, and applying the product rule to convection, proves the force estimate. Since $2\beta>1$, its square is integrable. Both clocks are increasing and unbounded, so Lemma~\ref{bcs:clock} applies to the event.
\end{proof}

The displayed bound applies to each mixed derivative separately; no claim of a common constant over all derivative orders is intended. An unslowed periodic field and any of these slowed fields define separate choices of force.

\begin{proposition}[A spatial field with a startup profile]\label{bcs:autonomous-profile}
Let $Z(x)$ be a smooth divergence-free field on a flat torus. Let $\sigma$ be the smooth step which is zero for $t\le0$, one for $t\ge1$, and satisfies $\sigma(t)+\sigma(1-t)=1$. Define
\[
 \alpha_0(t)=\sigma(2t-1),\qquad
 \alpha_1(t)=\frac{\sigma(2t-1)}{1+t},\qquad
 u(t,x)=\alpha(t)Z(x).
\]
For either choice $\alpha=\alpha_0,\alpha_1$, the residual force is
\begin{equation}\label{bcs:autonomous-force}
 f=\alpha'Z+\alpha^2(Z\cdot\nabla)Z-\nu\alpha\Delta Z.
\end{equation}
It has bounded mixed derivatives, gives zero initial velocity and pressure zero, and follows the spatial flow of $Z$ at the onto time $s(t)=\int_0^t\alpha$. For $\alpha_0$, $s(t)=t-3/4$ for $t\ge1$, and the force is stationary there. For $\alpha_1$,
\[
 s(t)=s(1)+\log\frac{1+t}{2}\quad(t\ge1),
\]
and the force tends uniformly to zero and belongs to space-time $L^2$. If $\int Z=0$, both forces have zero mean. In general their means equal $\alpha'\int Z$.
\end{proposition}
\begin{proof}
Substitution proves the force identity and the material correspondence. Symmetry of the step gives
$\int_0^1\sigma(2t-1)\,dt=\frac12\int_0^1\sigma(v)\,dv=1/4$.
The formulas for $s$ follow and show that its range is $[0,\infty)$. For the slow profile the coefficients $\alpha',\alpha^2,\alpha$ are respectively $O((1+t)^{-2})$, $O((1+t)^{-2})$, and $O((1+t)^{-1})$ after startup; all derivatives are bounded, and integration proves the force norm. Spatial integration of \eqref{bcs:autonomous-force} proves the mean statement.
\end{proof}

These elementary clock estimates do not give faster-than-every-power decay for an arbitrary repeated processor. For example, a nonnegative speed bounded by $C(1+t)^{-2}$ has finite accumulated time. The constructions with rapidly shrinking signals require their own memory mechanisms.
